\documentclass[11pt]{article}
\usepackage[margin=1.1in]{geometry}
\usepackage{amsmath, amssymb, amsthm, mathtools}
\usepackage{enumitem}
\usepackage{xcolor}
\usepackage[colorlinks=true, linkcolor=blue, citecolor=blue]{hyperref}

\definecolor{jobc}{RGB}{20,90,160}
\definecolor{havec}{RGB}{20,120,60}
\definecolor{needc}{RGB}{170,90,10}
\definecolor{riskc}{RGB}{170,30,30}

\newcommand{\job}[1]{\par\smallskip\noindent
  \textcolor{jobc}{\textbf{Job.}}\ \emph{#1}\par\smallskip}
\newcommand{\says}[1]{\par\noindent\textbf{Key claims.}\ #1\par\smallskip}
\newcommand{\have}[1]{\par\noindent
  \textcolor{havec}{\textbf{Have.}}\ #1\par}
\newcommand{\need}[1]{\par\noindent
  \textcolor{needc}{\textbf{Need.}}\ #1\par}
\newcommand{\risk}[1]{\par\noindent
  \textcolor{riskc}{\textbf{Risk.}}\ #1\par\smallskip}
\newcommand{\budget}[1]{\hfill{\small\textsf{[#1]}}}
\setlength{\emergencystretch}{3em}  % working doc: allow looser lines

% --- theorem environments and notation, kept identical to game_class.tex ---
\newtheorem{theorem}{Theorem}
\newtheorem{lemma}{Lemma}
\newtheorem{proposition}{Proposition}
\theoremstyle{definition}
\newtheorem{definition}{Definition}
\theoremstyle{remark}
\newtheorem{remark}{Remark}

\newcommand{\R}{\mathbb{R}}
\newcommand{\E}{\mathbb{E}}
\newcommand{\cA}{\mathcal{A}}
\newcommand{\cB}{\mathcal{B}}
\newcommand{\cP}{\mathcal{P}}
\newcommand{\cN}{\mathcal{N}}
\DeclareMathOperator*{\argmax}{arg\,max}
\DeclareMathOperator{\supp}{supp}

\title{Paper outline\\
\large Gaussian-mixture policies for continuous-action zero-sum games:\\
where they converge, and why}
\author{\texttt{cont\_tinkering} / \texttt{theory2}}
\date{}

\begin{document}
\maketitle

\begin{center}
\fbox{\parbox{0.9\textwidth}{\centering
\textbf{The final message --- everything builds to this line:}\\[4pt]
\emph{The failures practitioners attribute to tuning are properties of the game,
and they are checkable before you train.}}}
\end{center}

\bigskip
\noindent
\textbf{Spine.} The paper turns on one cheap lemma (\S\ref{sec:decomp}): the
mixture policy's update \emph{factors} into two subproblems --- \emph{where are
the atoms} and \emph{how much mass on each} --- and the second is a finite matrix
game solved globally and unconditionally by existing theory. Therefore all
difficulty, and all hypotheses, localise in the first. Every later section is a
consequence: \S\ref{sec:fail} shows how the first subproblem fails,
\S\ref{sec:class} characterises when it does not, \S\ref{sec:check} makes the
characterisation checkable, \S\ref{sec:valid} shows it is quantitatively
predictive, \S\ref{sec:poker} shows it survives contact with a real domain.

\medskip
\noindent
\textbf{Legend.} \textcolor{jobc}{\textbf{Job}} = what the section must
accomplish. \textcolor{havec}{\textbf{Have}} = material that already exists
(\texttt{theory2/game\_class.tex}, \texttt{check\_scale\_law.py}).
\textcolor{needc}{\textbf{Need}} = new work.
\textcolor{riskc}{\textbf{Risk}} = what a referee will attack.
Page budgets assume a 9-page main body.

\tableofcontents

\section*{Draft abstract}
\job{Written first, because it forces the message. Rewrite last.}

\noindent
Continuous action spaces in imperfect-information games --- bet sizing, bidding,
allocation --- are almost always handled by discretisation, at an abstraction cost
that is hand-tuned and unquantified. A natural alternative is a parametric
continuous policy, a Gaussian mixture, trained by a game-solving update. This
converges in practice, but no theory covers it: the convergence guarantees of
magnetic mirror descent require a monotone game, and in mixture \emph{parameters}
the game is badly non-monotone. We resolve this by observing that the update
factors exactly: the categorical head over components plays a finite matrix game,
which mirror descent solves globally and unconditionally, while the Gaussian
heads perform preconditioned ascent on a smoothed effective landscape. All
obstructions therefore live in the second problem --- \emph{transport} of the
component means to the equilibrium atoms --- and we characterise it. We show any
self-consistent pair of spurious local maxima is an exponentially stable non-Nash
attractor with an open basin, immune to the magnet, to extra capacity and to
annealing; this forces the shape of any initialisation-robust hypothesis. We then
give a class of games --- defined by a uniform one-point attraction condition on
the family of smoothed effective landscapes --- on which transport succeeds, show
the class is checkable in closed form for finite-rank couplings, and derive
design rules for the variance floor, the component count and the magnet. The
theory is quantitatively predictive: it predicts the residual exploitability
decays as $\Theta(1/t)$ and that the magnet rescales the whole Gaussian-head flow
by a closed-form factor, both confirmed to three or four digits. Finally we apply
the method to poker, where the near-finite support of equilibrium bet-size
distributions makes the class's central assumption a property of the domain
rather than a modelling convenience.

\section{Introduction}\label{sec:intro}
\budget{1.25 pages}

\job{Establish a \emph{puzzle}, not a gap. The reader should finish \S1 wanting
to know why an algorithm outside its own theory works at all --- and why it
sometimes doesn't.}

\begin{itemize}[leftmargin=1.2em]
\item Continuous actions in imperfect-information games: bet sizing in poker,
  bids in auctions, allocations in Blotto/security games. Standard practice
  discretises: hand-chosen grids, lossy, abstraction error unquantified.
\item The alternative: a Gaussian-mixture policy driven by a game-solving update.
  Empirically it converges.
\item \textbf{The puzzle.} Magnetic mirror descent guarantees convergence for
  \emph{monotone} games. In mixture parameters the game is wildly non-monotone
  (multi-well landscape). Yet the last iterate converges --- from some
  initialisations, and not others, with no apparent pattern. Practitioners
  respond by tuning std schedules, magnet coefficients, capacity, step sizes,
  largely without success.
\item \textbf{Thesis.} The erratic behaviour is neither noise nor tuning. The
  problem factors into two subproblems; one is unconditionally solved; the other
  fails for reasons that are a property of the \emph{game} and computable in
  advance.
\item Contribution list, in delivery order: (i) the factorisation; (ii) the
  sharpness result forcing the hypothesis; (iii) the class and its checkable
  instances; (iv) quantitative predictions verified to $10^{-3}$; (v) poker.
\end{itemize}

\risk{``Both ingredients already exist.'' Pre-empt: the contribution is the
\emph{characterisation}, not the algorithm. Say so in the first paragraph of
contributions, not in rebuttal.}

\section{Background}\label{sec:bg}
\budget{0.5 pages}

\job{Fix the game-theoretic object and the solution concept once, in exactly the
generality the rest of the paper uses --- compact action sets, bounded measurable
payoff, mixed strategies as measures --- so that \S\ref{sec:decomp} onwards can
speak of atoms, effective landscapes and exploitability without re-deriving
anything. Textbook material; stated because ``Nash equilibrium'' means something
weaker here than in the finite case and the reader must not import the finite
intuition.}

\subsection{The game}

\begin{definition}[Two-player zero-sum game with continuous action sets]
\label{def:game}
A \emph{two-player zero-sum game with continuous action sets} is a triple
$G = (\cA, \cB, u)$ in which
\begin{enumerate}[label=(\roman*), leftmargin=2.2em]
\item $\cA \subset \R^{n_0}$ and $\cB \subset \R^{n_1}$ are nonempty compact sets
  --- throughout this paper, compact convex bodies --- called the \emph{action
  sets} of player $0$ and player $1$;
\item $u : \cA \times \cB \to \R$ is bounded and Borel measurable, the
  \emph{payoff} (to player~$0$).
\end{enumerate}
The two players choose $a \in \cA$ and $b \in \cB$ simultaneously and
independently; player~$0$ then receives $u(a,b)$ and player~$1$ receives
$-u(a,b)$, so the payoffs sum to zero at every outcome. Player~$0$ maximises and
player~$1$ minimises. The game is \emph{continuous-action} in that $\cA$ and
$\cB$ are uncountable subsets of Euclidean space rather than finite sets; note
that \emph{no} continuity of $u$ in either argument is required.
\end{definition}

\begin{remark}[Asymmetry]\label{rem:asym}
Nothing ties the two sides together: $n_0 \ne n_1$, $\cA \ne \cB$, and later all
per-player hyperparameters (component counts, scale floors, step sizes, magnet
coefficients) may differ. Every statement in the paper is asymmetric in the
players unless it explicitly says otherwise.
\end{remark}

\paragraph{Mixed strategies.}
With $\cA$ and $\cB$ infinite there is in general no equilibrium in pure
actions, so we pass to the \emph{mixed extension}. Let $\cP(\cA)$ denote the set
of Borel probability measures on $\cA$; a \emph{mixed strategy} of player $0$ is
a $\mu \in \cP(\cA)$, and a pure action $a$ is identified with the Dirac measure
$\delta_a$. \emph{No topology on $\cP(\cA)$ is required here} --- nor anywhere up
to and including \eqref{eq:expl}; everything through the definition of Nash
equilibrium and of exploitability is purely measure-theoretic, and we introduce
the weak-$*$ topology only at the one point where it is actually used
(Remark~\ref{rem:weakstar}). The \emph{expected payoff} is
\begin{equation}\label{eq:U}
  U(\mu,\nu) \;\coloneqq\; \E_{\mu \otimes \nu}\,[\,u\,]
  \;=\; \int_\cA\!\!\int_\cB u(a,b)\,d\nu(b)\,d\mu(a),
  \qquad (\mu,\nu) \in \cP(\cA) \times \cP(\cB).
\end{equation}
Boundedness and measurability of $u$ make \eqref{eq:U} finite for every pair and,
by Fubini--Tonelli, independent of the order of integration; $U$ is
\emph{bilinear} (affine in each argument separately) and $|U| \le \|u\|_\infty$.
This is the only role the regularity assumptions on $u$ play: the difficulty of
the paper is never in defining the game.

\paragraph{Effective landscapes and pure best responses.}
Define the \emph{effective landscapes}
\begin{equation}\label{eq:eff}
  Q_\nu(a) \coloneqq \int_\cB u(a,b)\,d\nu(b),
  \qquad
  W_\mu(b) \coloneqq -\int_\cA u(a,b)\,d\mu(a),
\end{equation}
so that $U(\mu,\nu) = \int Q_\nu \,d\mu = -\int W_\mu \,d\nu$ and both players
\emph{maximise}: player $0$ maximises $\E_\mu Q_\nu$, player $1$ maximises
$\E_\nu W_\mu$. Because $U$ is affine in each argument, optimising over measures
is the same as optimising over points,
\begin{equation}\label{eq:pureBR}
  \sup_{\mu \in \cP(\cA)} U(\mu,\nu) \;=\; \sup_{a \in \cA} Q_\nu(a),
  \qquad
  \inf_{\nu \in \cP(\cB)} U(\mu,\nu) \;=\; -\sup_{b \in \cB} W_\mu(b),
\end{equation}
i.e.\ \emph{best responses may always be taken pure}. (The suprema are attained
when $Q_\nu$, resp.\ $W_\mu$, is upper semicontinuous --- in particular when $u$
is continuous; in general they are suprema.) Equation~\eqref{eq:pureBR} is what
lets a condition on a \emph{single} landscape $Q_\nu$ speak about the whole
mixed-strategy problem, and it is used silently throughout
\S\ref{sec:class}--\S\ref{sec:check}.

\subsection{Nash equilibrium}

\begin{definition}[Nash equilibrium]\label{def:nash}
A pair $(\mu^\star,\nu^\star) \in \cP(\cA) \times \cP(\cB)$ is a \emph{Nash
equilibrium} of $G$ if it is a saddle point of $U$:
\begin{equation}\label{eq:nash}
  U(\mu,\nu^\star) \;\le\; U(\mu^\star,\nu^\star) \;\le\; U(\mu^\star,\nu)
  \qquad \text{for all } \mu \in \cP(\cA),\ \nu \in \cP(\cB).
\end{equation}
Equivalently, neither player can strictly improve their own expected payoff by a
unilateral deviation: $\mu^\star$ maximises $U(\cdot,\nu^\star)$ and $\nu^\star$
minimises $U(\mu^\star,\cdot)$. By \eqref{eq:pureBR} it suffices to rule out
\emph{pure} deviations, so \eqref{eq:nash} is equivalent to
\[
  \sup_{a \in \cA} Q_{\nu^\star}(a)
  \;=\; U(\mu^\star,\nu^\star)
  \;=\; -\sup_{b \in \cB} W_{\mu^\star}(b),
\]
each player's landscape being flat at the equilibrium value across the other's
best responses.
\end{definition}

\begin{remark}[Value and interchangeability]\label{rem:value}
Weak duality $\sup_\mu \inf_\nu U \le \inf_\nu \sup_\mu U$ always holds; the
existence of a pair satisfying \eqref{eq:nash} forces equality, and the common
number
\[
  V(G) \;=\; \max_{\mu} \min_{\nu} U(\mu,\nu)
        \;=\; \min_{\nu} \max_{\mu} U(\mu,\nu)
        \;=\; U(\mu^\star,\nu^\star)
\]
is the \emph{value} of $G$. As in the finite case, the equilibrium set is a
product set: if $(\mu^\star,\nu^\star)$ and $(\mu',\nu')$ are equilibria then so
are $(\mu^\star,\nu')$ and $(\mu',\nu^\star)$, and all four have value $V(G)$.
Hence ``converging to a Nash equilibrium'' is a statement about the two marginals
separately --- which is what makes the asymmetric, per-player conditions of
\S\ref{sec:class} legitimate.
\end{remark}

\paragraph{Standing assumption: finitely supported equilibria.}
We fix one Nash equilibrium $(\mu^\star,\nu^\star)$ and assume throughout --- as
a hypothesis, not a consequence --- that both measures are purely atomic with
finite support,
\begin{equation}\label{eq:finite}
  \mu^\star = \sum_{j=1}^{M_0} w^\star_j\,\delta_{p_j},
  \qquad
  \nu^\star = \sum_{l=1}^{M_1} v^\star_l\,\delta_{r_l},
  \qquad
  w^\star \in \Delta_{M_0},\ v^\star \in \Delta_{M_1},\ w^\star,v^\star > 0 .
\end{equation}
The distinct points $p_j \in \cA$ and $r_l \in \cB$ are the \emph{atoms} of the
equilibrium. This is the assumption a $K$-component mixture policy needs in order
to be able to represent $\mu^\star$ at all, and it splits the learning problem in
two: finding the atom \emph{locations} is the job of the component means, finding
the \emph{weights} $w^\star, v^\star$ is the job of the categorical head. That
split is the factorisation of \S\ref{sec:decomp}. \S\ref{sec:poker} argues the
assumption is a property of the poker domain rather than a convenience.

\begin{remark}[Existence is assumed, not derived]\label{rem:exist}
If $u$ is \emph{continuous} on the compact set $\cA \times \cB$, a mixed Nash
equilibrium exists and the value exists (Ville; Glicksberg's extension of Nash's
theorem to compact Hausdorff action spaces), by a fixed-point argument over
$\cP(\cA) \times \cP(\cB)$ that needs the topology of
Remark~\ref{rem:weakstar}. For merely bounded measurable $u$ this can fail. We
therefore treat the existence of $(\mu^\star,\nu^\star)$ --- and separately its
finite support, \eqref{eq:finite} --- as standing assumptions rather than as
consequences. Cite this honestly: the paper is about \emph{reaching} an
equilibrium, not about its existence.
\end{remark}

\begin{remark}[Where a topology is needed --- and only there]\label{rem:weakstar}
Definitions~\ref{def:game} and~\ref{def:nash}, the reduction \eqref{eq:pureBR}
and the exploitability \eqref{eq:expl} are all measure-theoretic: they quantify
over measures and never over limits of measures. Topology enters at exactly
three later points, so we fix it once, here, for all of them. Equip $\cP(\cA)$
with the \emph{weak-$*$} topology --- the coarsest topology making
$\mu \mapsto \int f\,d\mu$ continuous for every $f \in C(\cA)$, equivalently
$\mu_k \to \mu$ iff $\int f\,d\mu_k \to \int f\,d\mu$ for all continuous $f$
(what probabilists call \emph{weak convergence}; the clash of names is worth a
footnote in the final text). As $\cA$ is compact metric, $\cP(\cA)$ is convex,
weak-$*$ compact and metrizable. Its three consumers:
\begin{enumerate}[label=(\roman*), leftmargin=2.2em, itemsep=1pt]
\item Remark~\ref{rem:exist}: Glicksberg's existence theorem needs exactly this
  compactness, together with continuity of $u$;
\item \S\ref{sec:valid}: releasing the variance floor gives
  $\cN(m,\sigma^2) \to \delta_m$ as $\sigma \downarrow 0$, a weak-$*$ statement
  that is \emph{false} in total variation, where
  $\|\delta_a - \delta_{a'}\|_{\mathrm{TV}} = 2$ whenever $a \ne a'$;
\item the branch-regularity condition (D) of \texttt{game\_class.tex}, whose
  ``metric metrising weak convergence'' is metrizability of this topology.
\end{enumerate}
Nothing else in the paper needs $\cP(\cA)$ to be topologised, and in particular
the class of \S\ref{sec:class} is stated without reference to it.
\end{remark}

\begin{remark}[Support characterisation]\label{rem:bg-support}
Applying \eqref{eq:nash} to $\mu = \delta_a$ gives $Q_{\nu^\star}(a) \le V(G)$
for all $a$, while $\int Q_{\nu^\star}\,d\mu^\star = V(G)$; hence
$Q_{\nu^\star} = V(G)$ $\mu^\star$-almost everywhere, and when $Q_{\nu^\star}$ is
upper semicontinuous,
\[
  \supp \mu^\star \subseteq \argmax_{\cA} Q_{\nu^\star},
  \qquad
  \supp \nu^\star \subseteq \argmax_{\cB} W_{\mu^\star}.
\]
\textbf{The equilibrium atoms sit at maxima of the effective landscape is
therefore a theorem, not a hypothesis} --- this is the sentence \S\ref{sec:class}
leans on when it assumes only nondegeneracy of those maxima and the absence of
others.
\end{remark}

\paragraph{Approximate equilibria: the metric we actually report.}
Convergence claims in \S\ref{sec:valid} are stated in \emph{exploitability}
(\textsc{NashConv}),
\begin{equation}\label{eq:expl}
  e(\mu,\nu) \;\coloneqq\;
  \sup_{\mu' \in \cP(\cA)} U(\mu',\nu) \;-\; \inf_{\nu' \in \cP(\cB)} U(\mu,\nu')
  \;\overset{\eqref{eq:pureBR}}{=}\;
  \sup_{a \in \cA} Q_\nu(a) \;+\; \sup_{b \in \cB} W_\mu(b) \;\ge\; 0 .
\end{equation}
A pair is an \emph{$\varepsilon$-Nash equilibrium} if $e(\mu,\nu) \le
\varepsilon$, and $e(\mu,\nu) = 0$ exactly when $(\mu,\nu)$ is a Nash
equilibrium in the sense of Definition~\ref{def:nash}. Note that $e$ is
computable from the two effective landscapes alone, which is why it is the
natural error measure for a paper whose whole analysis is landscape-based, and it
is the quantity predicted to decay as $\Theta(1/t)$ in \S\ref{sec:valid}.

\have{\texttt{game\_class.tex} \S1.1: the payoff, the effective landscapes
\eqref{eq:eff}, the finite-support assumption \eqref{eq:finite} and
Remark~\ref{rem:bg-support} are already written and can be lifted verbatim.}

\need{Nothing mathematically new. Compress hard --- this is background and must
not read as contribution. If space is tight, Definitions~\ref{def:game}
and~\ref{def:nash} plus \eqref{eq:pureBR} and \eqref{eq:expl} stay in the main
body; Remarks~\ref{rem:value}--\ref{rem:exist} move to the appendix.}

\risk{A referee reads a half-page of standard definitions as padding. Mitigate by
pointing forward in every remark: each one is here because a later section uses
it (\eqref{eq:pureBR} $\to$ \S\ref{sec:class}, Remark~\ref{rem:value} $\to$
asymmetric conditions, Remark~\ref{rem:bg-support} $\to$ the ``atoms are maxima''
step).}

\section{Setting and algorithm}\label{sec:setup}
\budget{0.75 pages}

\job{Define the object precisely and get out of the way. No results here.}

\says{Two-player zero-sum game on compact action sets; $K$-component
diagonal-Gaussian mixture policy; the two-head update (closed-form entropic
simplex step on the categorical head, Fisher--Rao natural gradient on
means/log-stds, magnet snapshot refreshed every $T$ steps); variance floor and
ceiling.}

\have{\texttt{game\_class.tex} \S1.1--1.2 verbatim.}

\need{One paragraph positioning the algorithm as \textbf{SAC $\cap$ MMD}:
reparameterised Gaussian policy and entropy from the former, KL-proximal magnet
from the latter (the magnet is a target network by another name). Two sentences
of closest related work; full treatment deferred to \S\ref{sec:related}.}

\risk{Do \emph{not} headline ``SAC + MMD''. It invites ``so it is a combination
of two known things''. The algorithm is the setting; the theory is the paper.}

\section{The decomposition}\label{sec:decomp}
\budget{1 page --- the pivot of the paper}

\job{The reader should feel the problem change shape. After this section the
question is no longer ``when does the algorithm converge'' but ``when do the
means find the atoms''.}

\says{With $\bar A_{kl}$ the expected payoff of component $k$ against component
$l$, the expected payoff is $U = w^\top \bar A\, v$, and hence:
\emph{(i)} the categorical heads run \emph{exactly} tabular MMD on the finite
matrix game $\bar A$; \emph{(ii)} the Gaussian heads run preconditioned ascent on
the smoothed effective landscape at each component's own scale. Requires no
structure on the payoff whatsoever.}

\begin{itemize}[leftmargin=1.2em]
\item \textbf{Corollary.} Matrix-game MMD converges last-iterate from any
  initialisation --- no basins, no assumptions. So the weight subproblem is
  \emph{solved}, and every obstruction lives in the means.
\item \textbf{Remark to include here}: this is why a \emph{one-point} condition
  can characterise a \emph{many-atom} equilibrium. The target of the later
  condition is an \emph{atom}, not an equilibrium; the mixing is not geometric,
  it is a matrix game. Readers arriving from the optimisation literature will
  otherwise stumble on exactly this.
\end{itemize}

\have{Lemma~1 and the ``target is an atom'' remark, both proved.}

\section{When transport fails}\label{sec:fail}
\budget{1 page}

\job{Make the hypothesis of \S\ref{sec:class} feel \emph{forced} rather than
chosen. Negative result before positive result --- this is the rhetorical hinge of
the paper.}

\says{Any self-consistent pair of spurious local maxima of the two effective
landscapes is a fixed point of the dynamics, and is locally exponentially stable
--- hence has an \emph{open basin} of initialisations from which the method does
not converge to Nash.}

\begin{itemize}[leftmargin=1.2em]
\item It is \textbf{magnet-invariant}: at the trap the magnet is refreshed to the
  trap itself, so its gradient vanishes; and it adds $-\eta\tau I$ to the
  Jacobian, strictly \emph{deepening} the trap.
\item It is \textbf{capacity-invariant}: spare components are captured too.
\item It is \textbf{annealing-invariant} --- worse, annealing is the mechanism
  that delivers components into it.
\item \textbf{Therefore} any initialisation-robust class must forbid spurious
  maxima outright. The shape of condition (C) is derived, not chosen.
\item Worked counterexample: the decoy game, excluded by mode counting alone
  (three local maxima, support size two) --- no dominance computation needed.
\end{itemize}

\have{Both propositions, proved, plus the counterexample.}

\section{The class}\label{sec:class}
\budget{1.75 pages}

\job{State the class minimally and prove what can be proved. Resist completeness:
eight conditions in the body will lose the reader.}

\says{Core in the body --- \textbf{(C)} cellwise star-concavity of the smoothed,
magnet-regularised effective landscape, uniformly over opponent policies and over
the scale range; \textbf{(D)} the branch maps are Lipschitz with round-trip gain
below one; \textbf{(M)} the branch endpoints at equilibrium are the atoms;
\textbf{(K)} capacity $K_i \ge M_i$ with an initialisation meeting every cell.}

\begin{itemize}[leftmargin=1.2em]
\item Push the remaining conditions (nondegeneracy, atom regularity, step-size,
  weight non-degeneracy) to an appendix as regularity and bookkeeping --- all four
  turned out weaker than they first appeared.
\item Emphasise what (C) is \emph{not}: not concavity (a Gaussian bump is concave
  only within its own width, yet convergence happens from far outside), and not
  ``the basin contains the target'' (basins of a time-varying flow are not
  invariant). Uniformity in the opponent is the load-bearing part.
\item (D) is the genuinely \emph{game-theoretic} condition --- it has no
  counterpart in single-agent graduated optimisation, because there is no
  opponent to chase.
\item Main result plus the proof architecture: three input-to-state-stable blocks
  (means, scales, weights) composed by a small-gain theorem.
\item \textbf{State plainly what (K) does not do}: the class does not solve
  exploration. Deterministic local ascent cannot move a component between
  basins, so covering is assumed because the algorithm genuinely cannot achieve
  it. Better to own this than to be caught by it.
\end{itemize}

\have{Definition, the ISS architecture, the scale lemma, the local results.}

\need{\textbf{The global theorem is not proved} --- see
\S\ref{sec:decisions}, decision 1. This determines the shape of the section.}

\risk{The strongest attack on the paper. Do not paper over it.}

\section{Checkable instances and design rules}\label{sec:check}
\budget{1 page}

\job{Convert the class into things a practitioner does differently on Monday.
This is what makes the theory useful rather than decorative.}

\begin{itemize}[leftmargin=1.2em]
\item \textbf{Finite-rank coupling}: the opponent enters through $m$ moments, so
  ``for all $\nu$'' collapses to a quantifier over a compact convex moment body.
  Curvature-dominates-coupling in closed form, with worked constants.
\item \textbf{One-dimensional Gaussian-bump landscapes}: scale-space causality
  means a single mode count at the floor plus a merge check certifies the whole
  scale ladder. Structure finer than the floor is free.
\item \textbf{Multiple dimensions}: what survives verbatim (everything
  structural) and what degrades (effective mass $h w^M$, smoothing can
  \emph{create} maxima, covering costs more components, bias gains a factor $M$).
\item \textbf{Design rules that fall out}: floor below half the atom spacing;
  $K \gg M$ with a spread initialisation; magnet only on flat directions;
  $\varepsilon = \Theta(1/t)$ converts a target exploitability into a step budget.
\end{itemize}

\have{All of it.}

\section{Quantitative validation}\label{sec:valid}
\budget{1.25 pages}

\job{Show the theory is \emph{predictive}, not merely consistent. Parameter-free
numerical predictions are rare in this literature; lead with that framing.}

\begin{itemize}[leftmargin=1.2em]
\item Scale decay $\sigma \sim t^{-1/2}$ and residual exploitability
  $\varepsilon = \Theta(1/t)$: fitted constant vs.\ predicted, four digits;
  $\sigma^{-2}$ affine in $t$ with $R^2 \ge 1-3\times10^{-8}$. State the mode of
  convergence when the floor is released: the policy tends to $\mu^\star$
  \emph{weak-$*$} (Remark~\ref{rem:weakstar}(ii)), which is the only sense in
  which a Gaussian mixture can approach an atomic measure at all --- the
  exploitability rate is the quantitative form of that statement.
\item The magnet acts as a pure \textbf{time-rescaling} of the Gaussian-head flow
  by a closed-form factor; verified over a $16\times$ range, with configurations
  sharing $\eta\tau T$ collapsing onto one constant.
\item Overcapacity \emph{helps} inside the class and \emph{fails} outside it ---
  the same knob with opposite signs, predicted in advance.
\item The certified counterexample, and the observation that residuals reported
  in prior work were snapshots of a decaying quantity rather than floor-set
  constants.
\end{itemize}

\have{\texttt{check\_scale\_law.py}, runs A--F, results in
\texttt{results/scale\_law.\{json,log\}}.}

\need{Figures: (a) $\sigma^{-2}$ vs.\ $t$, three runs, one line; (b) brake factor
vs.\ $\eta\tau T$, measured vs.\ closed form; (c) basin maps inside vs.\ outside
the class.}

\section{Poker}\label{sec:poker}
\budget{1.75 pages}

\job{Show the class survives contact with a real domain, in three rungs of
increasing realism and decreasing rigour --- and label the rungs honestly.}

\subsection*{8.1 Why poker fits}
\begin{itemize}[leftmargin=1.2em]
\item The action space is \emph{hybrid} --- discrete (fold/check/call/bet) plus a
  continuous bet size --- which is exactly the shape of the policy. No contortion.
\item \textbf{Near-finite support of equilibrium bet-size distributions is a known
  property of the domain}: solvers discretise to a handful of sizes and lose
  little. The class's central assumption is therefore domain-justified, not a
  modelling convenience. This is the single best argument for the pairing.
\item EV as a function of bet size is kinked and can jump (calling thresholds
  move discretely). The class needs only \emph{bounded measurable} payoffs,
  because the variance floor does the regularising --- a place where this theory
  is better suited to poker than a smoothness-based one would be.
\item Capacity: the number of sizes per node is unknown and varies; $K \gg M$ is
  safe inside the class, so the practitioner does not need to know $M$.
\end{itemize}

\subsection*{8.2 Analytic poker models --- the bridge}
\job{The rung that makes 8.3 interpretable. Do not skip it.}
\begin{itemize}[leftmargin=1.2em]
\item Von Neumann poker / one-street continuous-bet games: genuinely poker,
  genuinely continuous, \emph{closed-form equilibria}.
\item Verify (C) numerically on the true EV-vs-size landscape; recover the known
  optimal sizes; measure \textbf{exact} exploitability.
\item Report the mode structure of the EV landscape --- this is where a spurious
  maximum would show up if the polarised-range intuition creates one.
\end{itemize}
\need{Everything. Model implementation, landscape diagnostics, exploitability.}

\subsection*{8.3 Heads-up no-limit hold'em --- the scaling demonstration}
\begin{itemize}[leftmargin=1.2em]
\item Continuous bet sizing \emph{without discretisation}; baseline is a matched
  discretised-sizing agent.
\item Evaluation: local best response, plus head-to-head against the baseline.
\item \textbf{Labelled explicitly as outside the theorem's hypotheses.}
\end{itemize}
\need{Environment, network, self-play infrastructure, evaluation.}
\risk{Extensive form is a second theory, not a corollary: payoffs are bilinear in
sequence form, but the network parameterises \emph{behavioural} strategies per
infoset, and the cellwise analysis would have to be redone with continuation
values as the landscape. Say this in \S\ref{sec:limits} before a referee does.}

\section{Related work}\label{sec:related}
\budget{0.5 pages}

\job{Position honestly on four axes.}
\begin{itemize}[leftmargin=1.2em]
\item Game solving with mirror descent / magnetic mirror descent; last-iterate
  results in monotone and bilinear games.
\item Continuous and separable games; finite-support equilibria.
\item Graduated optimisation and $\sigma$-niceness --- the single-agent ancestor of
  condition (C); state precisely what the game setting adds (the uniformity in
  the opponent, and the small-gain condition, which have no single-agent
  counterpart).
\item Star-convexity / quasar-convexity in nonconvex optimisation, and
  variational stability in learning-in-games --- the two literatures the
  inequality is borrowed from. Be explicit that the novelty is the
  \emph{conjunction}: cellwise on a partition, toward a moving target, uniform in
  the opponent, uniform in scale.
\item Action abstraction in poker; whether continuous bet-size learning has been
  attempted.
\end{itemize}
\need{A literature check on the last point before claiming novelty.}

\section{Limitations}\label{sec:limits}
\budget{0.25 pages}

\job{Disarm the obvious attacks by making them first. Each one costs a sentence
here and would cost the paper a referee otherwise.}
\begin{itemize}[leftmargin=1.2em]
\item Extensive form is a separate theory.
\item (C) is unverifiable at scale; checkability holds for finite-rank couplings
  and low-dimensional landscapes.
\item Exploration is assumed, not solved --- the class inherits the covering
  condition because the algorithm cannot discharge it.
\item Shared network parameters couple infosets; the analysis treats mixture
  parameters as free.
\item Exact gradients; the stochastic case is open.
\end{itemize}

\section{Conclusion}\label{sec:conc}
\budget{0.25 pages}

\job{One paragraph. Return to the final message and make it land.}

\noindent
The two subproblems are not equally hard: equilibration is free, transport is
everything. Transport succeeds or fails according to the geometry of the smoothed
effective landscape, which is a property of the game --- computable before
training, not discoverable by sweeping hyperparameters. \emph{The failures
practitioners attribute to tuning are properties of the game, and they are
checkable before you train.}

\appendix
\section{Deferred conditions and proofs}
Nondegeneracy, atom regularity, step-size and weight conditions; full proofs;
the resolution lemma and its reading as a hyperparameter design rule; the
per-direction magnet taxonomy.

\section{Open decisions}\label{sec:decisions}

\job{Settle these \emph{before} drafting --- both change \S\ref{sec:intro}'s
contribution list.}

\paragraph{Decision 1: the unproved global theorem.}
Proved today: the decomposition, the sharpness proposition, local convergence,
the scale and bias laws. Not proved: the small-gain composition of the three ISS
blocks. Three options, ranked:
\begin{enumerate}[leftmargin=1.6em]
\item \textbf{Close it under a clean extra hypothesis} --- a two-timescale limit,
  or an explicit small-gain constant. Yields a genuine theorem; narrows the class
  slightly. \emph{Preferred.}
\item \textbf{Restructure} so the theorems are the negative and structural
  results, with the global statement as a criterion supported by the ISS
  architecture and the experiments. Honest, publishable, weaker headline.
\item Claim it as proved. \emph{No.}
\end{enumerate}

\paragraph{Decision 2: rung 8.2.}
The analytic poker models are the only place where poker and verifiable theory
meet. Without them, \S\ref{sec:poker} is a leap from synthetic games to an
unverifiable benchmark. Recommendation: build 8.2 even at the cost of shrinking
8.3.

\paragraph{Decision 3: venue and framing.}
A theory-first framing (class + sharpness, poker as demonstration) and an
empirics-first framing (poker headline, theory as justification) want different
venues and different \S\ref{sec:intro}s. The material is currently much stronger
as theory-first.

\end{document}
